\introduction
L'éducation joue un rôle fondamental dans le développement d'un pays, 
et la gestion des événements interactifs, des cours et des formations
 constitue un aspect clé pour assurer une expérience d'apprentissage 
 efficace et enrichissante. Cependant, malgré la présence de plateformes 
 existantes comme Slido, qui offrent des fonctionnalités interactives, 
 il demeure un besoin crucial pour une solution mobile intégrée et 
 personnalisée. En effet, les outils actuels ne répondent pas pleinement 
 aux exigences des organisateurs et des participants, notamment en matière
  d'accessibilité, d'interactivité fluide et de gestion optimale des événements. 
  Cette situation complique la tâche des formateurs et des participants, 
  limitant ainsi l'efficacité des événements et des sessions de formation.
  C'est dans ce cadre que notre projet prend forme, visant à développer 
  une application mobile dédiée à la gestion 
  d'événements interactifs, inspirée de Slido, mais en offrant des 
  fonctionnalités supplémentaires adaptées aux 
  besoins spécifiques des organisateurs et des participants.
   Cette solution numérique vise à enrichir l'expérience 
  des utilisateurs en facilitant les interactions en temps 
  réel, en offrant une accessibilité maximale, et en permettant 
  la gestion des événements de manière fluide et intuitive.

\section*{Problématique}

Dans le cadre des événements interactifs et des formations actuelles, il existe 
un manque d'intégration et de personnalisation des outils utilisés.
 Les solutions existantes, bien que fonctionnelles, ne parviennent 
 pas à offrir une gestion optimale des interactions en temps réel et 
à garantir l'accessibilité pour tous les participants, en particulier 
dans des contextes où la connectivité peut être un obstacle. Cela
engendre une gestion fragmentée des événements, une communication
peu fluide et une limitation des retours d'information instantanés, 
nécessaires pour améliorer l'efficacité des événements. Il est donc 
impératif de proposer une solution intégrée et adaptée, capable de 
répondre à ces défis et de stimuler un engagement optimal.

\section*{Contexte et Justification}
\cite{ World-Bank}
Face à ces enjeux, notre projet consiste à développer une application  \gls{TIC} 
mobile interactive, permettant de gérer les événements, les cours et les 
formations de manière efficace. Cette solution intégrera des fonctionnalités 
de Slido pour faciliter les interactions en temps réel (Q\&A, sondages, quiz) 
tout en offrant des outils supplémentaires, tels que la gestion des participants
 et des statistiques détaillées. L'application proposera également un mode hors 
 ligne, essentiel pour garantir l'accessibilité dans des contextes où la connexion
 Internet est limitée. En améliorant l'accessibilité et la communication, ce projet 
 ambitionne de transformer l'organisation des événements et de maximiser l'impact
 des formations et des cours.

\section*{Objectif}

L'objectif principal de ce projet est de développer une application mobile 
interactive et inclusive, qui vise à améliorer la gestion des événements et
 à faciliter les interactions en temps réel. Plus précisément, il s'agira de :

\begin{itemize}
    \item \textbf{Améliorer l'accessibilité et la personnalisation} : Offrir une 
    interface adaptable aux besoins spécifiques des organisateurs, avec une 
    gestion fine des rôles et permissions.
    \item \textbf{Faciliter les interactions en temps réel} : Intégrer et optimiser
     les fonctionnalités interactives (sondages, quiz) pour une meilleure 
     visualisation des résultats et une interaction fluide.
    \item \textbf{Proposer un mode hors ligne} : Développer un système permettant 
    aux utilisateurs d'interagir et de consulter les supports, même sans connexion 
    Internet.
    \item \textbf{Optimiser la gestion des participants} : Suivre l'engagement 
    des participants en temps réel, avec des statistiques détaillées pour
     améliorer les événements futurs.
\end{itemize}

\glsresetall

%  \gls{acro} puis \Gls{acroglo} et enfin \gls{glossaire}


